\section{光滑单位分解}
\label{sec:mfd-partition-unity}

很多构造在一张坐标图里很容易完成，例如在坐标中写下一个函数、一个向量场或一个内积。但整个流形通常需要多张图，各图上的构造在重叠处一般不一致，不能直接拼接。办法是给每张图配一个光滑权重，只在该图内部取正值，所有权重之和处处为一；各图上的构造乘以权重后相加，在每点附近只涉及有限项，结果就是整体上光滑的对象。本节构造这样的权重。

\paragraph{支撑受限的光滑函数}
第一步要有在一个区域外恒为零、在区域内为正的光滑函数。多项式做不到这一点：非零多项式不能在一个区间上恒为零。下面的函数在原点处以所有阶导数为零的方式“起飞”：
\begin{equation}
\eta(t)=\begin{cases}e^{-1/t},&t>0,\\0,&t\leq0.\end{cases}
\label{eq:mfd-smooth-step}
\end{equation}
当$t>0$时，$\eta$的$k$阶导数形如$P_k(1/t)e^{-1/t}$，$P_k$是多项式；令$s=1/t\to+\infty$，$P_k(s)e^{-s}\to0$，所以各阶导数从右侧趋于零，与左侧的零函数衔接，$\eta$在$\mathbb R$上光滑。

用$\eta$可以造出欧氏空间中的光滑截断函数。对$0<r<R$，令
\[
\chi(\bm x)=\frac{\eta(R^2-\lVert\bm x\rVert_2^2)}{\eta(R^2-\lVert\bm x\rVert_2^2)+\eta(\lVert\bm x\rVert_2^2-r^2)} .
\]
分母的两项不会同时为零，因为$\lVert\bm x\rVert_2^2<R^2$与$\lVert\bm x\rVert_2^2>r^2$至少有一个成立。于是$\chi$光滑，取值在$[0,1]$中，在半径$r$的闭球上恒为$1$，在半径$R$的球外恒为$0$。这类函数称为\term{截断函数}（cutoff function）。

函数$f$的\term{支撑}（support）是集合$\{p:f(p)\neq0\}$的闭包，记为$\operatorname{supp}f$。取闭包是为了让“$f$在支撑之外恒为零”成为开集上的性质：$\chi$的支撑是半径$R$的闭球。把$\chi$放进一张坐标图，并乘到坐标表示上，就能把局部构造截断成在整个流形上有定义的函数，前提是截断后的支撑落在图的定义域内部。

\paragraph{局部有限的覆盖}
多个权重相加时，若每点附近有无穷多个权重非零，和可能不收敛，也可能不光滑。所以需要覆盖满足局部有限条件。设$\{V_\beta\}$是$M$的一族子集，若每点都有一个邻域只与其中有限个$V_\beta$相交，则称这族子集\term{局部有限}（locally finite）。若一个开覆盖$\{V_\beta\}$的每个成员都含于给定开覆盖$\{U_\alpha\}$的某个成员中，则称前者是后者的一个\term{细化}（refinement）。

\begin{definition}[光滑单位分解]
\label{def:mfd-partition-unity}
设$\{U_\alpha\}_{\alpha\in A}$是$M$的开覆盖。一族光滑函数$\{\psi_\alpha\}_{\alpha\in A}$称为从属于$\{U_\alpha\}$的光滑\term{单位分解}（partition of unity），若：
\begin{enumerate}
\item 对每个$\alpha$，$0\leq\psi_\alpha\leq1$且$\operatorname{supp}\psi_\alpha\subseteq U_\alpha$；
\item 支撑族$\{\operatorname{supp}\psi_\alpha\}$局部有限；
\item 对每个$p\in M$，$\sum_{\alpha}\psi_\alpha(p)=1$。
\end{enumerate}
\end{definition}

由局部有限性，每点附近只有有限个$\psi_\alpha$不为零，所以第三条中的和在每点附近是有限和，自然光滑。

\paragraph{一维的例子}
在$\mathbb R$上取开覆盖$U_1=(-\infty,1)$、$U_2=(-1,\infty)$。令
\[
\psi_1(x)=\frac{\eta(\tfrac12-x)}{\eta(\tfrac12-x)+\eta(x+\tfrac12)},\qquad
\psi_2(x)=\frac{\eta(x+\tfrac12)}{\eta(\tfrac12-x)+\eta(x+\tfrac12)} .
\]
对每个$x$，$x<\tfrac12$与$x>-\tfrac12$至少有一个成立，分母为正。两函数都光滑，取值于$[0,1]$，和为$1$。$\psi_1$的支撑是$(-\infty,\tfrac12]$，含于$U_1$；$\psi_2$的支撑是$[-\tfrac12,\infty)$，含于$U_2$。

这里特意把截断点取在$\pm\tfrac12$而不是$\pm1$。若用$\eta(1-x)$，支撑就是$(-\infty,1]$，不含于开集$U_1$。支撑必须严格落在开集内，乘上权重的局部构造才能在开集的边缘以外光滑地补零。

\paragraph{存在定理}
单位分解的存在要用到流形定义中的全部条件：局部欧氏性提供坐标球上的截断函数，Hausdorff性与第二可数性保证能把覆盖细化为局部有限的。

\begin{theorem}[光滑单位分解的存在]
\label{thm:mfd-partition-unity-existence}
设$M$是光滑流形，可以带边界，$\{U_\alpha\}_{\alpha\in A}$是$M$的任一开覆盖。则存在从属于$\{U_\alpha\}$的光滑单位分解。
\end{theorem}

\begin{proof}
先处理$M$紧的情形，它已包含构造的全部想法。对每个$p\in M$，取包含$p$的坐标图$(V,\varphi)$与某个$U_\alpha\ni p$，在$\varphi(V\cap U_\alpha)$中取以$\varphi(p)$为中心、半径为$R_p$的球，其闭包仍在其中。令$B_p$为半径$R_p/2$的同心球在$\varphi^{-1}$下的原像。边界点处用半球代替球，下同。

$\{B_p\}$覆盖$M$，由紧性可选出有限个$B_{p_1},\ldots,B_{p_m}$覆盖$M$。对每个$k$，取截断函数$\chi_k$，在半径$R_{p_k}/2$的闭球上为$1$，在半径$\tfrac34R_{p_k}$的球外为$0$。令$\lambda_k=\chi_k\circ\varphi_k$在图内，图外为$0$。$\lambda_k$的支撑是半径$\tfrac34R_{p_k}$的闭球的原像，它是紧集，在Hausdorff空间中是闭集，所以图外补零后$\lambda_k$在$M$上光滑。

$\lambda_k\geq0$，并在$B_{p_k}$上等于$1$，所以$\Lambda=\sum_k\lambda_k$处处至少为$1$。令$\tilde\psi_k=\lambda_k/\Lambda$，它们光滑、非负、和为一，且每个$\tilde\psi_k$的支撑含于某个$U_{\alpha(k)}$。最后对每个$\alpha$令$\psi_\alpha=\sum_{k:\alpha(k)=\alpha}\tilde\psi_k$，空和取$0$，就得到以原指标集编号的单位分解。

一般情形只需把“选有限个$B_p$”换成“选局部有限的可数个”。第二可数性与局部欧氏性给出一列紧集$K_1\subseteq K_2\subseteq\cdots$，满足$K_j$含于$K_{j+1}$的内部且$\bigcup_jK_j=M$。对每个$j$，紧集$K_{j+1}\setminus\operatorname{int}K_j$可用有限个如上的坐标球覆盖，并让这些球含于$\operatorname{int}K_{j+2}\setminus K_{j-1}$。每点只与相邻几层的球相交，所以全部球组成局部有限的覆盖，其余步骤与紧情形相同。紧集列的构造这里不展开，它是第二可数局部紧Hausdorff空间的一般性质。
\end{proof}

\paragraph{单位分解的用途}
单位分解的典型用法是“局部构造，加权求和”。设$f_\alpha$是定义在$U_\alpha$上的光滑函数，则$\sum_\alpha\psi_\alpha f_\alpha$在整个$M$上有定义并且光滑，其中每项在$U_\alpha$外按零理解。若所有$f_\alpha$在某点附近都等于同一个值，加权和在那里也等于这个值，因为权重之和为一。

由此可得两个直接结论。第一，任一点$p$的任一邻域中都有光滑函数在$p$附近等于某个给定的局部光滑函数，在邻域外为零；这说明局部定义的坐标函数可以截断成全局光滑函数，切向量完全由它对全局光滑函数的作用确定。第二，在每张图中取欧氏内积，用单位分解加权相加，就在每个切空间上得到一个内积。正定性在加权求和下保持，因为每点处的权重非负、和为一，至少有一项严格为正。这正是第~\ref{chap:riemannian-geometry}章黎曼度量存在的依据。同样的手法也用于把局部坐标中的积分拼成流形上的积分，见第~\ref{chap:manifold-calculus}章。
